\documentclass[aspectratio=169,11pt]{ctexbeamer}

\usetheme{Warsaw}
\setbeamertemplate{navigation symbols}{}

\usepackage{booktabs}
\usepackage{tabularx}
\usepackage{fancyvrb}
\usepackage{xcolor}

\definecolor{codebg}{RGB}{248,248,248}
\setbeamercolor{block body}{bg=codebg,fg=black}

\title{Python 绘图工具与命令概览}
\subtitle{JupyterLab、NumPy、Matplotlib、Plotly 与 ipywidgets}
\author{多元微积分可视化示例配套讲义}
\date{\today}

\begin{document}

\begin{frame}
  \titlepage
\end{frame}

\begin{frame}{本讲定位}
  \begin{itemize}
    \item 介绍 notebook 中出现的绘图工具和命令的一般使用方法。
    \item 重点不是某一个 notebook 的操作流程，而是可迁移到其他课程、实验和数据分析任务的绘图模式。
    \item 主要覆盖：JupyterLab、NumPy 网格、Matplotlib 二维图、Plotly 交互图、ipywidgets 参数交互。
  \end{itemize}
\end{frame}

\begin{frame}{典型 Python 绘图流程}
  \begin{enumerate}
    \item 准备数据：数组、表格、网格或函数采样点。
    \item 创建图形对象：Matplotlib 的 \texttt{fig, ax}，或 Plotly 的 \texttt{fig}。
    \item 添加图层：曲线、散点、等高线、向量场、曲面。
    \item 调整样式：标题、坐标轴、颜色、比例、图例、色条。
    \item 显示或导出：在 Jupyter 中显示，或保存为图片/PDF/HTML。
  \end{enumerate}
\end{frame}

\begin{frame}[fragile]{JupyterLab：交互式绘图环境}
  JupyterLab 适合教学和探索式计算，因为文字、公式、代码和图形可以放在同一个文档中。

  \begin{block}{常用启动方式}
\begin{Verbatim}[fontsize=\small]
jupyter lab
jupyter lab some_notebook.ipynb
\end{Verbatim}
  \end{block}

  \begin{itemize}
    \item Code cell 执行 Python，Markdown cell 写解释和公式。
    \item Kernel 保存变量和函数；重启 kernel 后需要重新执行定义代码。
    \item 图形输出直接显示在 cell 下方，适合边改参数边观察结果。
  \end{itemize}
\end{frame}

\begin{frame}[fragile]{依赖导入：约定俗成的别名}
  大多数科学计算代码会使用固定别名，便于阅读和交流。

  \begin{block}{常见导入}
\begin{Verbatim}[fontsize=\small]
import numpy as np
import matplotlib.pyplot as plt
import plotly.graph_objects as go
from plotly.subplots import make_subplots
from ipywidgets import FloatSlider, interact
\end{Verbatim}
  \end{block}

  \begin{itemize}
    \item \texttt{np}：数组、网格、线性代数和数学函数。
    \item \texttt{plt}：Matplotlib 的状态式入口，常配合面向对象的 \texttt{ax} 使用。
    \item \texttt{go}：Plotly 的图层对象，例如 \texttt{go.Surface}、\texttt{go.Contour}。
  \end{itemize}
\end{frame}

\begin{frame}[fragile]{NumPy：一维坐标轴}
  绘图前通常先生成坐标点。连续区间需要被离散成有限个采样点。

  \begin{block}{等距采样}
\begin{Verbatim}[fontsize=\small]
x = np.linspace(-3, 3, 121)
y = np.linspace(-3, 3, 121)
\end{Verbatim}
  \end{block}

  \begin{itemize}
    \item \texttt{np.linspace(a, b, n)} 在闭区间 \([a,b]\) 上生成 \(n\) 个等距点。
    \item 点数越多，图形越平滑，但计算量和图形体积也越大。
    \item 一维坐标轴可用于曲线图，也可继续扩展成二维网格。
  \end{itemize}
\end{frame}

\begin{frame}[fragile]{NumPy：二维网格}
  二元函数 \(z=f(x,y)\) 的绘图一般需要二维网格。

  \begin{block}{网格生成与函数采样}
\begin{Verbatim}[fontsize=\small]
X, Y = np.meshgrid(x, y)
Z = f(X, Y)
\end{Verbatim}
  \end{block}

  \begin{itemize}
    \item \texttt{X} 和 \texttt{Y} 是同形状二维数组，表示每个网格点的坐标。
    \item \texttt{f(X,Y)} 要支持 NumPy 数组输入，称为向量化计算。
    \item \texttt{Z} 与 \texttt{X,Y} 同形状，可直接用于等高线和曲面图。
  \end{itemize}
\end{frame}

\begin{frame}[fragile]{Matplotlib：Figure 与 Axes}
  Matplotlib 推荐使用面向对象接口：\texttt{fig} 表示整张图，\texttt{ax} 表示一个坐标区域。

  \begin{block}{基础模板}
\begin{Verbatim}[fontsize=\small]
fig, ax = plt.subplots(figsize=(7, 5))
ax.plot(x, y)
ax.set_title("标题")
ax.set_xlabel("x")
ax.set_ylabel("y")
ax.grid(alpha=0.2)
plt.show()
\end{Verbatim}
  \end{block}
\end{frame}

\begin{frame}[fragile]{Matplotlib：等高线 \texttt{contour}}
  等高线用于展示二元函数 \(f(x,y)=c\) 的水平集。

  \begin{block}{等高线命令}
\begin{Verbatim}[fontsize=\small]
contour = ax.contour(X, Y, Z, levels=16, cmap="viridis")
ax.clabel(contour, inline=True, fontsize=8)
ax.set_aspect("equal", adjustable="box")
\end{Verbatim}
  \end{block}

  \begin{itemize}
    \item \texttt{levels} 控制等高线数量或指定具体函数值。
    \item \texttt{cmap} 控制颜色映射。
    \item \texttt{ax.clabel} 给等高线添加数值标签。
    \item \texttt{set\_aspect("equal")} 保持横纵坐标单位长度一致。
  \end{itemize}
\end{frame}

\begin{frame}[fragile]{Matplotlib：填色等高线 \texttt{contourf}}
  \texttt{contour} 画线，\texttt{contourf} 给等高线区域填色。

  \begin{block}{区域填色与色条}
\begin{Verbatim}[fontsize=\small]
filled = ax.contourf(X, Y, Z, levels=24, cmap="viridis")
fig.colorbar(filled, ax=ax, shrink=0.85)
\end{Verbatim}
  \end{block}

  \begin{itemize}
    \item \texttt{contourf} 更适合展示函数值整体分布。
    \item \texttt{fig.colorbar} 说明颜色与数值的对应关系。
    \item 线状等高线和填色等高线也可以叠加使用。
  \end{itemize}
\end{frame}

\begin{frame}[fragile]{Matplotlib：向量场 \texttt{quiver}}
  \texttt{quiver} 用箭头展示二维向量场，例如梯度场。

  \begin{block}{向量场命令}
\begin{Verbatim}[fontsize=\small]
quiver = ax.quiver(QX, QY, U, V, C,
                   cmap="magma", pivot="mid", scale=28)
fig.colorbar(quiver, ax=ax)
\end{Verbatim}
  \end{block}

  \begin{itemize}
    \item \texttt{QX,QY} 是箭头起点坐标。
    \item \texttt{U,V} 是箭头的水平和竖直分量。
    \item \texttt{C} 可选，用颜色表达另一个数值，例如向量长度。
    \item \texttt{pivot="mid"} 表示箭头以中点锚定。
  \end{itemize}
\end{frame}

\begin{frame}[fragile]{梯度箭头：方向与长度分开}
  梯度很大时，直接画原始向量会导致箭头过长。常见做法是方向归一化，长度用颜色表示。

  \begin{block}{归一化方向}
\begin{Verbatim}[fontsize=\small]
speed = np.hypot(GX, GY)
U = np.divide(GX, speed, out=np.zeros_like(GX),
              where=speed > 1e-12)
V = np.divide(GY, speed, out=np.zeros_like(GY),
              where=speed > 1e-12)
ax.quiver(QX, QY, U, V, speed, cmap="magma")
\end{Verbatim}
  \end{block}

  \texttt{where=speed > 1e-12} 用来避免在零向量处除以零。
\end{frame}

\begin{frame}[fragile]{Matplotlib：标注与图例}
  标注命令让图形更适合讲解和发布。

  \begin{block}{常见标注命令}
\begin{Verbatim}[fontsize=\small]
ax.scatter(px, py, s=45, c="black", label="选定点")
ax.annotate("(1.5, 1.0)", (1.5, 1.0),
            xytext=(6, 6), textcoords="offset points")
ax.legend(loc="upper right")
\end{Verbatim}
  \end{block}

  \begin{itemize}
    \item \texttt{scatter} 标出离散点。
    \item \texttt{annotate} 添加文本说明。
    \item 设置 \texttt{label} 后，再调用 \texttt{legend} 才会显示图例。
  \end{itemize}
\end{frame}

\begin{frame}[fragile]{Matplotlib：中文字体}
  Matplotlib 默认字体不一定包含中文字形，所以中文标题、图例、色条可能显示为方框。

  \begin{block}{常用配置}
\begin{Verbatim}[fontsize=\small]
plt.rcParams["font.family"] = "sans-serif"
plt.rcParams["font.sans-serif"] = ["Noto Sans CJK SC",
                                   "Microsoft YaHei",
                                   "SimHei",
                                   "DejaVu Sans"]
plt.rcParams["axes.unicode_minus"] = False
\end{Verbatim}
  \end{block}

  Linux/WSL 可安装 \texttt{fonts-noto-cjk}，必要时删除 \texttt{\textasciitilde/.cache/matplotlib} 后重启 kernel。
\end{frame}

\begin{frame}[fragile]{Plotly：Figure 与 Trace}
  Plotly 的核心结构是 \texttt{Figure + Trace}：一张图可以包含多个图层。

  \begin{block}{基本模式}
\begin{Verbatim}[fontsize=\small]
fig = go.Figure()
fig.add_trace(go.Scatter(x=x, y=y, mode="lines"))
fig.update_layout(title="交互式图形")
fig.show()
\end{Verbatim}
  \end{block}

  \begin{itemize}
    \item \texttt{Trace} 是具体图层，例如曲线、散点、曲面、等高线。
    \item \texttt{update\_layout} 调整整张图的标题、尺寸、边距等。
    \item 在 JupyterLab 中，Plotly 图通常可以直接缩放、拖动和悬停查看数值。
  \end{itemize}
\end{frame}

\begin{frame}[fragile]{Plotly：三维曲面 \texttt{go.Surface}}
  三维曲面适合展示二元函数 \(z=f(x,y)\) 的整体形状。

  \begin{block}{曲面图层}
\begin{Verbatim}[fontsize=\small]
surface = go.Surface(
    x=X,
    y=Y,
    z=Z,
    colorscale="Viridis",
    contours={"z": {"show": True, "usecolormap": True}},
)
fig.add_trace(surface)
\end{Verbatim}
  \end{block}

  \begin{itemize}
    \item \texttt{x,y,z} 可以都是二维数组。
    \item \texttt{colorscale} 用颜色表达高度或函数值。
    \item \texttt{contours} 可在曲面上叠加等高线痕迹。
  \end{itemize}
\end{frame}

\begin{frame}[fragile]{Plotly：等高线 \texttt{go.Contour}}
  Plotly 也可以绘制可交互的二维等高线图。

  \begin{block}{等高线图层}
\begin{Verbatim}[fontsize=\small]
contour = go.Contour(
    x=x,
    y=y,
    z=Z,
    colorscale="Viridis",
    contours={"showlabels": True,
              "labelfont": {"size": 10}},
)
fig.add_trace(contour)
\end{Verbatim}
  \end{block}

  \texttt{x,y} 通常是一维坐标轴，\texttt{z} 是二维函数值矩阵。
\end{frame}

\begin{frame}[fragile]{Plotly：子图 \texttt{make\_subplots}}
  多图并排时，可以用 \texttt{make\_subplots} 组织布局。

  \begin{block}{混合三维曲面和二维等高线}
\begin{Verbatim}[fontsize=\small]
fig = make_subplots(
    rows=1,
    cols=2,
    specs=[[{"type": "surface"}, {"type": "contour"}]],
)
fig.add_trace(go.Surface(x=X, y=Y, z=Z), row=1, col=1)
fig.add_trace(go.Contour(x=x, y=y, z=Z), row=1, col=2)
\end{Verbatim}
  \end{block}

  \texttt{specs} 必须声明每个子图类型，否则三维图层可能无法放入对应位置。
\end{frame}

\begin{frame}[fragile]{Plotly：布局与场景}
  Plotly 的样式通常通过 \texttt{update\_layout}、\texttt{update\_xaxes}、\texttt{update\_scenes} 等方法设置。

  \begin{block}{常用设置}
\begin{Verbatim}[fontsize=\small]
fig.update_layout(width=1000, height=560, title="标题")
fig.update_scenes(
    xaxis_title="x",
    yaxis_title="y",
    zaxis_title="z",
    aspectmode="cube",
)
fig.update_xaxes(title_text="x", row=1, col=2)
\end{Verbatim}
  \end{block}
\end{frame}

\begin{frame}[fragile]{ipywidgets：把参数变成滑块}
  \texttt{interact} 会把控件的值传给函数，并在值变化时重新运行函数。

  \begin{block}{滑块模板}
\begin{Verbatim}[fontsize=\small]
def draw(a=1.0, h=0.0, d=1.0):
    A = np.array([[a, h], [h, d]])
    # 计算并绘图

interact(
    draw,
    a=FloatSlider(value=1.0, min=-2, max=2, step=0.25),
    h=FloatSlider(value=0.0, min=-2, max=2, step=0.25),
    d=FloatSlider(value=1.0, min=-2, max=2, step=0.25),
)
\end{Verbatim}
  \end{block}
\end{frame}

\section{Notebook 示例代码}

\begin{frame}[fragile]{示例代码：曲面与等高线（一）}
  下面是 notebook 中用于生成曲面和等高线的三个二元函数示例。

  \begin{block}{碗形曲面与鞍面}
\begin{Verbatim}[fontsize=\tiny]
def bowl(x, y):
    return x**2 + y**2

plot_surface_and_contours(
    bowl,
    xlim=(-3, 3),
    ylim=(-3, 3),
    title="例 1：f(x, y) = x² + y²",
)

def saddle(x, y):
    return x**2 - y**2

plot_surface_and_contours(
    saddle,
    xlim=(-3, 3),
    ylim=(-3, 3),
    title="例 2：f(x, y) = x² - y²",
)
\end{Verbatim}
  \end{block}
\end{frame}

\begin{frame}[fragile]{示例代码：曲面与等高线（二）}
  周期函数示例使用更大的观察区域，常用 \([-2\pi,2\pi]^2\)。

  \begin{block}{周期函数曲面}
\begin{Verbatim}[fontsize=\tiny]
def wave(x, y):
    return np.sin(x) * np.cos(y)

plot_surface_and_contours(
    wave,
    xlim=(-2 * np.pi, 2 * np.pi),
    ylim=(-2 * np.pi, 2 * np.pi),
    title="例 3：f(x, y) = sin(x)cos(y)",
)
\end{Verbatim}
  \end{block}

  这个例子体现了 \texttt{np.sin}、\texttt{np.cos} 对数组输入的向量化支持。
\end{frame}

\begin{frame}[fragile]{示例代码：梯度与等高线（一）}
  梯度函数和原函数分开定义，再一起传给绘图函数。

  \begin{block}{\(f(x,y)=x^2+y^2\) 的梯度}
\begin{Verbatim}[fontsize=\tiny]
def grad_bowl(x, y):
    return 2 * x, 2 * y

plot_contour_with_gradient(
    bowl,
    grad_bowl,
    xlim=(-3, 3),
    ylim=(-3, 3),
    points=[(1.5, 1.0), (-2.0, 1.0)],
    title="f(x, y) = x² + y² 的等高线与梯度方向",
)
\end{Verbatim}
  \end{block}
\end{frame}

\begin{frame}[fragile]{示例代码：梯度与等高线（二）}
  在指定点画出梯度方向和等高线切向，可以直接观察二者正交。

  \begin{block}{单点方向比较}
\begin{Verbatim}[fontsize=\tiny]
plot_gradient_and_tangent(
    bowl,
    grad_bowl,
    point=(1.5, 1.0),
    xlim=(-3, 3),
    ylim=(-3, 3),
    title="在一点比较梯度方向和等高线切向",
)
\end{Verbatim}
  \end{block}
\end{frame}

\begin{frame}[fragile]{示例代码：梯度与等高线（三）}
  鞍面示例中，梯度在 \(x\) 方向和 \(y\) 方向的符号不同。

  \begin{block}{\(f(x,y)=x^2-y^2\) 的梯度}
\begin{Verbatim}[fontsize=\tiny]
def grad_saddle(x, y):
    return 2 * x, -2 * y

plot_contour_with_gradient(
    saddle,
    grad_saddle,
    xlim=(-3, 3),
    ylim=(-3, 3),
    points=[(1.5, 1.0), (1.5, -1.0)],
    title="f(x, y) = x² - y² 的等高线与梯度方向",
)
\end{Verbatim}
  \end{block}
\end{frame}

\begin{frame}[fragile]{示例代码：二次型曲面分类（一）}
  二次型示例用矩阵 \(A\) 表示二次项，并交给统一绘图函数处理。

  \begin{block}{正定、负定、不定、半正定}
\begin{Verbatim}[fontsize=\tiny]
A_positive = np.array([[1, 0], [0, 1]])
plot_quadratic_surface(
    A_positive,
    title="正定矩阵：f(x, y) = x² + y²",
)

A_negative = np.array([[-1, 0], [0, -1]])
plot_quadratic_surface(
    A_negative,
    title="负定矩阵：f(x, y) = -x² - y²",
)

A_indefinite = np.array([[1, 0], [0, -1]])
plot_quadratic_surface(
    A_indefinite,
    title="不定矩阵：f(x, y) = x² - y²",
)
\end{Verbatim}
  \end{block}
\end{frame}

\begin{frame}[fragile]{示例代码：二次型曲面分类（二）}
  半正定和含线性项的例子展示了退化与平移。

  \begin{block}{半正定与线性项}
\begin{Verbatim}[fontsize=\tiny]
A_semipositive = np.array([[1, 0], [0, 0]])
plot_quadratic_surface(
    A_semipositive,
    title="半正定矩阵：f(x, y) = x²",
)

plot_quadratic_surface(
    A_positive,
    b=(-2, 4),
    c=0,
    xlim=(-3, 5),
    ylim=(-6, 2),
    title="含线性项：f(x, y) = x² + y² - 2x + 4y",
)
\end{Verbatim}
  \end{block}
\end{frame}

\begin{frame}[fragile]{示例代码：交互式二次曲面（一）}
  交互函数负责根据滑块参数构造矩阵、分类并绘图。

  \begin{block}{交互函数}
\begin{Verbatim}[fontsize=\tiny]
def explore_quadratic(a=1.0, h=0.0, d=1.0):
    A = np.array([[a, h], [h, d]], dtype=float)
    info = classify_quadratic(A)
    eigen_text = np.array2string(
        info["eigenvalues"],
        precision=3,
    )
    print(f"A = [[{a:.2f}, {h:.2f}], [{h:.2f}, {d:.2f}]]")
    print(f"特征值：{eigen_text}")
    print(f"分类：{info['classification']}")
    f = quadratic_function(A)
    plot_surface_and_contours(
        f,
        xlim=(-3, 3),
        ylim=(-3, 3),
        n=81,
        title="交互式二次曲面",
    )
\end{Verbatim}
  \end{block}
\end{frame}

\begin{frame}[fragile]{示例代码：交互式二次曲面（二）}
  \texttt{interact} 把三个矩阵参数变成滑块。

  \begin{block}{滑块绑定}
\begin{Verbatim}[fontsize=\tiny]
interact(
    explore_quadratic,
    a=FloatSlider(
        value=1.0, min=-2.0, max=2.0,
        step=0.25, description="a",
    ),
    h=FloatSlider(
        value=0.0, min=-2.0, max=2.0,
        step=0.25, description="h",
    ),
    d=FloatSlider(
        value=1.0, min=-2.0, max=2.0,
        step=0.25, description="d",
    ),
);
\end{Verbatim}
  \end{block}
\end{frame}

\begin{frame}[fragile]{示例代码：练习模板}
  练习代码的设计思路是只让学习者改函数或矩阵，其他绘图细节保持封装。

  \begin{block}{自定义函数与矩阵}
\begin{Verbatim}[fontsize=\tiny]
def f_student(x, y):
    return (
        np.exp(-0.25 * (x**2 + y**2))
        * np.cos(2 * x)
        + 0.15 * y
    )

plot_surface_and_contours(
    f_student,
    xlim=(-4, 4),
    ylim=(-4, 4),
    title="练习 1：自定义函数",
)

A_student = np.array([[1.0, 0.75], [0.75, -0.5]])
plot_quadratic_surface(
    A_student,
    title="练习 2：自定义二次型矩阵",
)
\end{Verbatim}
  \end{block}
\end{frame}

\begin{frame}{何时用 Matplotlib，何时用 Plotly}
  \begin{tabularx}{\linewidth}{@{}lXX@{}}
    \toprule
    场景 & Matplotlib & Plotly \\
    \midrule
    课堂静态讲解 & 稳定，易标注，易导出 & 可用，但交互优势不明显 \\
    二维等高线/向量场 & 很适合 & 适合悬停和交互查看 \\
    三维曲面 & 可画，但交互弱 & 更适合旋转和缩放 \\
    论文/讲义图片 & 更可控 & 需额外导出配置 \\
    浏览器交互 & 一般 & 原生优势 \\
    \bottomrule
  \end{tabularx}
\end{frame}

\begin{frame}{常见错误与排查}
  \begin{tabularx}{\linewidth}{@{}lX@{}}
    \toprule
    现象 & 排查方向 \\
    \midrule
    \texttt{ModuleNotFoundError} & 依赖未安装，检查 Python 环境和 \texttt{pip install} 位置。 \\
    图中中文是方框 & 配置中文字体，安装 \texttt{fonts-noto-cjk} 或系统中文字体。 \\
    \texttt{Z.shape} 不对 & 函数没有正确支持 NumPy 数组输入，检查是否用了非向量化写法。 \\
    箭头太密或太乱 & 降低 \texttt{density}，归一化方向，或调 \texttt{scale}。 \\
    Plotly 子图报类型错误 & 检查 \texttt{specs} 中的子图类型是否匹配 trace。 \\
    \bottomrule
  \end{tabularx}
\end{frame}

\begin{frame}{实践建议}
  \begin{itemize}
    \item 先写清楚数据结构：一维数组、二维网格、还是矩阵。
    \item 优先使用 Matplotlib 的 \texttt{fig, ax = plt.subplots()} 模式。
    \item 二维解释图重视坐标比例、标签、图例和色条。
    \item 三维图优先用 Plotly，并设置合适的 \texttt{aspectmode} 和相机视角。
    \item 交互控件只控制少量关键参数，避免一次暴露太多选择。
  \end{itemize}
\end{frame}

\begin{frame}{总结}
  \begin{itemize}
    \item NumPy 负责把数学对象变成可计算、可绘制的数组。
    \item Matplotlib 负责可控的二维静态绘图。
    \item Plotly 负责浏览器中的交互式二维/三维绘图。
    \item ipywidgets 负责把参数探索变成交互控件。
    \item JupyterLab 把文字、公式、代码和图形整合成适合教学和探索的工作流。
  \end{itemize}
\end{frame}

\section{参数曲线与隐式几何绘图工具}

\begin{frame}{本专题新增的绘图对象}
  \begin{itemize}
    \item 参数曲线 \(r(t)=(x(t),y(t))\)：\texttt{np.linspace}、\texttt{ax.plot}、\texttt{np.gradient}、\texttt{ax.quiver}。
    \item 隐式曲线 \(F(x,y)=c\)：\texttt{np.meshgrid}、\texttt{ax.contour}、\texttt{ax.clabel}。
    \item 参数曲面 \(r(u,v)\)：\texttt{np.meshgrid} 配合 \texttt{indexing="ij"}，再用 \texttt{go.Surface}。
    \item 隐式曲面 \(F(x,y,z)=c\)：三维网格、\texttt{marching\_cubes}、\texttt{go.Mesh3d}。
    \item 参数实验：\texttt{FloatSlider}、\texttt{interact}。
  \end{itemize}
\end{frame}

\begin{frame}[fragile]{环境扩展：隐式曲面需要 scikit-image}
  二维等高线可直接由 Matplotlib 完成；三维隐式曲面通常先提取三角网格，再交给 Plotly 绘制。

  \begin{block}{notebook 中的核心导入}
\begin{Verbatim}[fontsize=\tiny]
import numpy as np
import matplotlib.pyplot as plt
import plotly.graph_objects as go
from ipywidgets import FloatSlider, interact
from skimage.measure import marching_cubes
\end{Verbatim}
  \end{block}

  \begin{itemize}
    \item \texttt{skimage.measure.marching\_cubes} 把 \(F(x,y,z)=c\) 的采样体数据转成三角面片。
    \item \texttt{plotly.graph\_objects.Mesh3d} 接收顶点和三角形索引，显示可旋转的曲面。
  \end{itemize}
\end{frame}

\begin{frame}[fragile]{二维参数曲线：先采样参数}
  参数曲线的基本思想是：先取一列参数 \(t\)，再计算对应的 \(x(t),y(t)\)。

  \begin{block}{单位圆例子}
\begin{Verbatim}[fontsize=\tiny]
t = np.linspace(0, 2 * np.pi, 600)
x_circle = np.cos(t)
y_circle = np.sin(t)

plot_parametric_curve(
    x_circle,
    y_circle,
    t,
    tangent_at=[np.pi / 6, 2 * np.pi / 3],
    title="参数曲线例 1：单位圆 r(t)=(cos t, sin t)",
)
\end{Verbatim}
  \end{block}

  \texttt{x\_circle}、\texttt{y\_circle}、\texttt{t} 是同长度一维数组。
\end{frame}

\begin{frame}[fragile]{Matplotlib 画参数曲线}
  一条参数曲线本质上仍是平面折线图，关键命令是 \texttt{ax.plot(x,y)}。

  \begin{block}{绘制曲线、起点和终点}
\begin{Verbatim}[fontsize=\tiny]
fig, ax = plt.subplots(figsize=(7.2, 6.2))
ax.plot(x, y, color="royalblue", linewidth=2.4, label="参数曲线")
ax.scatter([x[0]], [y[0]], color="green", s=45, label="起点", zorder=4)
ax.scatter([x[-1]], [y[-1]], color="black", s=45, label="终点", zorder=4)
ax.set_aspect("equal", adjustable="box")
ax.grid(alpha=0.25)
ax.legend(loc="best")
plt.show()
\end{Verbatim}
  \end{block}
\end{frame}

\begin{frame}[fragile]{参数曲线的切向量}
  没有解析导数时，可用 \texttt{np.gradient} 对采样点做数值微分。

  \begin{block}{从采样点计算并绘制切向量}
\begin{Verbatim}[fontsize=\tiny]
dx_dt = np.gradient(x, t)
dy_dt = np.gradient(y, t)

idx = int(np.argmin(np.abs(t - t0)))
tangent = normalize_vector([dx_dt[idx], dy_dt[idx]]) * arrow_scale
ax.quiver(
    [x[idx]], [y[idx]],
    [tangent[0]], [tangent[1]],
    angles="xy", scale_units="xy", scale=1,
    color="crimson", width=0.008,
)
\end{Verbatim}
  \end{block}

  \texttt{scale\_units="xy"} 表示箭头长度按坐标轴单位解释。
\end{frame}

\begin{frame}[fragile]{参数曲线的更多例子}
  改变 \(x(t),y(t)\) 的公式，就可以复用同一个绘图函数。

  \begin{block}{螺线与旋轮线}
\begin{Verbatim}[fontsize=\tiny]
t_spiral = np.linspace(0, 5 * np.pi, 900)
r = 0.08 * t_spiral
x_spiral = r * np.cos(t_spiral)
y_spiral = r * np.sin(t_spiral)

t_cycloid = np.linspace(0, 4 * np.pi, 900)
x_cycloid = t_cycloid - np.sin(t_cycloid)
y_cycloid = 1 - np.cos(t_cycloid)
\end{Verbatim}
  \end{block}

  同样的数据结构可以画闭曲线、开曲线、自交曲线和周期轨迹。
\end{frame}

\begin{frame}[fragile]{隐式曲线：二维网格采样}
  隐式曲线 \(F(x,y)=c\) 不能直接写成一条参数序列，通常先在平面网格上计算 \(F\)。

  \begin{block}{二维网格工具}
\begin{Verbatim}[fontsize=\tiny]
def make_2d_grid(xlim=(-3, 3), ylim=(-3, 3), n=401):
    x = np.linspace(xlim[0], xlim[1], n)
    y = np.linspace(ylim[0], ylim[1], n)
    X, Y = np.meshgrid(x, y)
    return x, y, X, Y

_, _, X, Y = make_2d_grid(xlim, ylim, n)
Z = np.asarray(F(X, Y), dtype=float)
\end{Verbatim}
  \end{block}

  \texttt{F} 需要支持数组输入，例如 \texttt{x**2 + y**2 - 1}。
\end{frame}

\begin{frame}[fragile]{Matplotlib 画隐式曲线}
  \texttt{ax.contour} 可以从网格函数值中提取水平集。

  \begin{block}{背景等高线与目标水平集}
\begin{Verbatim}[fontsize=\tiny]
background_levels = np.linspace(np.nanmin(Z), np.nanmax(Z), 18)
ax.contour(X, Y, Z, levels=background_levels,
           cmap="Greys", alpha=0.45, linewidths=0.7)

contour = ax.contour(
    X, Y, Z,
    levels=[level],
    colors="crimson",
    linewidths=2.5,
)
ax.clabel(contour, fmt={level: f"F={level:g}"}, inline=True, fontsize=9)
\end{Verbatim}
  \end{block}
\end{frame}

\begin{frame}[fragile]{隐式曲线的法向与切向}
  对 \(F(x,y)=c\)，梯度 \(\nabla F=(F_x,F_y)\) 垂直于水平集；二维切向可取 \((-F_y,F_x)\)。

  \begin{block}{法向与切向箭头}
\begin{Verbatim}[fontsize=\tiny]
gx, gy = grad_F(px, py) if grad_F else numerical_gradient_2d(F, px, py)
normal = normalize_vector([gx, gy]) * arrow_scale
tangent = normalize_vector([-gy, gx]) * arrow_scale

ax.quiver([px], [py], [normal[0]], [normal[1]],
          angles="xy", scale_units="xy", scale=1,
          color="crimson", width=0.008, label="法向 grad F")
ax.quiver([px], [py], [tangent[0]], [tangent[1]],
          angles="xy", scale_units="xy", scale=1,
          color="royalblue", width=0.008, label="切向")
\end{Verbatim}
  \end{block}
\end{frame}

\begin{frame}[fragile]{隐式曲线例子}
  notebook 中把隐式函数、梯度和采样范围分开写，便于替换。

  \begin{block}{圆与双曲线}
\begin{Verbatim}[fontsize=\tiny]
def F_circle(x, y):
    return x**2 + y**2 - 1

def grad_F_circle(x, y):
    return 2 * x, 2 * y

def F_hyperbola(x, y):
    return x**2 - y**2 - 1

def grad_F_hyperbola(x, y):
    return 2 * x, -2 * y
\end{Verbatim}
  \end{block}
\end{frame}

\begin{frame}[fragile]{三维参数曲面：两个参数轴}
  参数曲面需要两个参数 \(u,v\)。\texttt{indexing="ij"} 让数组轴顺序对应 \(u,v\)。

  \begin{block}{球面的参数采样}
\begin{Verbatim}[fontsize=\tiny]
u = np.linspace(0.05, np.pi - 0.05, 90)
v = np.linspace(0, 2 * np.pi, 120)
U, V = np.meshgrid(u, v, indexing="ij")

X_sphere = np.sin(U) * np.cos(V)
Y_sphere = np.sin(U) * np.sin(V)
Z_sphere = np.cos(U)
\end{Verbatim}
  \end{block}

  \texttt{X\_sphere,Y\_sphere,Z\_sphere} 是同形状二维数组。
\end{frame}

\begin{frame}[fragile]{Plotly 画参数曲面}
  Plotly 的 \texttt{go.Surface} 直接接收 \(X,Y,Z\) 三个二维数组。

  \begin{block}{参数曲面图层}
\begin{Verbatim}[fontsize=\tiny]
fig = go.Figure()
fig.add_trace(
    go.Surface(
        x=X,
        y=Y,
        z=Z,
        colorscale=colorscale,
        opacity=0.92,
        colorbar={"title": "参数"},
        name="参数曲面",
    )
)
fig.update_scenes(xaxis_title="x", yaxis_title="y",
                  zaxis_title="z", aspectmode="data")
fig.show()
\end{Verbatim}
  \end{block}
\end{frame}

\begin{frame}[fragile]{参数曲面的切平面信息}
  对 \(r(u,v)\)，两个偏导 \(r_u,r_v\) 是切向量，叉乘 \(r_u\times r_v\) 是法向量。

  \begin{block}{从采样曲面计算切向和法向}
\begin{Verbatim}[fontsize=\tiny]
dX_du, dX_dv = np.gradient(X, u, v, edge_order=2)
dY_du, dY_dv = np.gradient(Y, u, v, edge_order=2)
dZ_du, dZ_dv = np.gradient(Z, u, v, edge_order=2)

ru = np.array([dX_du[i, j], dY_du[i, j], dZ_du[i, j]])
rv = np.array([dX_dv[i, j], dY_dv[i, j], dZ_dv[i, j]])
normal = np.cross(ru, rv)
\end{Verbatim}
  \end{block}

  \texttt{np.cross} 是三维向量叉乘，常用于曲面法向量。
\end{frame}

\begin{frame}[fragile]{Plotly 画三维向量}
  notebook 用 \texttt{go.Scatter3d} 的线段和端点标记来表示三维箭头。

  \begin{block}{向已有 Figure 添加一条三维向量}
\begin{Verbatim}[fontsize=\tiny]
origin = np.asarray(origin, dtype=float)
vector = np.asarray(vector, dtype=float) * scale
end = origin + vector

fig.add_trace(
    go.Scatter3d(
        x=[origin[0], end[0]],
        y=[origin[1], end[1]],
        z=[origin[2], end[2]],
        mode="lines+markers",
        line={"color": color, "width": 7},
        marker={"size": 4, "color": color},
        name=name,
    )
)
\end{Verbatim}
  \end{block}
\end{frame}

\begin{frame}[fragile]{参数曲面例子}
  环面和螺旋面仍然遵循同一结构：生成 \(u,v\)，计算 \(X,Y,Z\)，调用绘图函数。

  \begin{block}{环面参数方程}
\begin{Verbatim}[fontsize=\tiny]
u = np.linspace(0, 2 * np.pi, 120)
v = np.linspace(0, 2 * np.pi, 80)
U, V = np.meshgrid(u, v, indexing="ij")

R, a = 2.0, 0.55
X_torus = (R + a * np.cos(V)) * np.cos(U)
Y_torus = (R + a * np.cos(V)) * np.sin(U)
Z_torus = a * np.sin(V)
\end{Verbatim}
  \end{block}
\end{frame}

\begin{frame}[fragile]{隐式曲面：三维网格采样}
  隐式曲面 \(F(x,y,z)=c\) 需要在三维盒子中采样函数值。

  \begin{block}{三维网格工具}
\begin{Verbatim}[fontsize=\tiny]
def make_3d_grid(xlim=(-2, 2), ylim=(-2, 2), zlim=(-2, 2), n=64):
    x = np.linspace(xlim[0], xlim[1], n)
    y = np.linspace(ylim[0], ylim[1], n)
    z = np.linspace(zlim[0], zlim[1], n)
    X, Y, Z = np.meshgrid(x, y, z, indexing="ij")
    return x, y, z, X, Y, Z

x, y, z, X, Y, Z = make_3d_grid(xlim, ylim, zlim, n)
values = np.asarray(F(X, Y, Z), dtype=float)
\end{Verbatim}
  \end{block}
\end{frame}

\begin{frame}[fragile]{marching cubes：从体数据到三角网格}
  \texttt{marching\_cubes} 返回顶点坐标和三角形顶点索引。

  \begin{block}{提取等值面}
\begin{Verbatim}[fontsize=\tiny]
dx = (xlim[1] - xlim[0]) / (n - 1)
dy = (ylim[1] - ylim[0]) / (n - 1)
dz = (zlim[1] - zlim[0]) / (n - 1)

vertices, faces, _, _ = marching_cubes(
    values,
    level=level,
    spacing=(dx, dy, dz),
)
vertices += np.array([xlim[0], ylim[0], zlim[0]])
\end{Verbatim}
  \end{block}

  \texttt{spacing} 和平移修正用于把数组索引坐标转换回真实 \(x,y,z\) 坐标。
\end{frame}

\begin{frame}[fragile]{Plotly Mesh3d：显示隐式曲面}
  \texttt{go.Mesh3d} 使用 \texttt{x,y,z} 存顶点，用 \texttt{i,j,k} 存三角形三个顶点的索引。

  \begin{block}{三角网格图层}
\begin{Verbatim}[fontsize=\tiny]
fig = go.Figure(
    data=[
        go.Mesh3d(
            x=vertices[:, 0], y=vertices[:, 1], z=vertices[:, 2],
            i=faces[:, 0], j=faces[:, 1], k=faces[:, 2],
            color="steelblue",
            opacity=0.62,
            flatshading=False,
            lighting={"ambient": 0.55, "diffuse": 0.75, "specular": 0.25},
            name="隐式曲面",
        )
    ]
)
\end{Verbatim}
  \end{block}
\end{frame}

\begin{frame}[fragile]{隐式曲面的法向量}
  对 \(F(x,y,z)=c\)，\(\nabla F=(F_x,F_y,F_z)\) 是曲面的法向量。

  \begin{block}{球面例子}
\begin{Verbatim}[fontsize=\tiny]
def F_sphere(x, y, z):
    return x**2 + y**2 + z**2 - 1

def grad_F_sphere(x, y, z):
    return 2 * x, 2 * y, 2 * z

p = 1 / np.sqrt(3)
plot_implicit_surface(
    F_sphere,
    xlim=(-1.4, 1.4), ylim=(-1.4, 1.4), zlim=(-1.4, 1.4),
    n=64,
    grad_F=grad_F_sphere,
    normal_point=(p, p, p),
    title="隐式曲面例 1：球面 x²+y²+z²-1=0，grad F 是法向量",
)
\end{Verbatim}
  \end{block}
\end{frame}

\begin{frame}[fragile]{隐式曲面例子：环面}
  隐式环面适合展示 \texttt{marching\_cubes + Mesh3d} 的价值，因为它不能写成 \(z=f(x,y)\) 的单值函数。

  \begin{block}{环面隐式方程}
\begin{Verbatim}[fontsize=\tiny]
def F_torus(x, y, z, R=1.55, r=0.45):
    return (
        (x**2 + y**2 + z**2 + R**2 - r**2) ** 2
        - 4 * R**2 * (x**2 + y**2)
    )

plot_implicit_surface(
    lambda x, y, z: F_torus(x, y, z, R=1.55, r=0.45),
    xlim=(-2.2, 2.2), ylim=(-2.2, 2.2), zlim=(-0.9, 0.9),
    n=72,
    title="隐式曲面例 4：环面",
)
\end{Verbatim}
  \end{block}
\end{frame}

\begin{frame}[fragile]{交互控件：水平集参数}
  \texttt{interact} 适合把一个常数或少量形状参数变成滑块。

  \begin{block}{圆的水平集}
\begin{Verbatim}[fontsize=\tiny]
def explore_circle_level(c=1.0):
    plot_implicit_curve(
        lambda x, y: x**2 + y**2,
        xlim=(-2.2, 2.2),
        ylim=(-2.2, 2.2),
        level=c,
        n=301,
        title=f"交互实验 1：水平集 x²+y²={c:.2f}",
    )

interact(
    explore_circle_level,
    c=FloatSlider(value=1.0, min=0.1, max=4.0, step=0.1, description="c"),
);
\end{Verbatim}
  \end{block}
\end{frame}

\begin{frame}[fragile]{交互控件：隐式曲面参数}
  三维交互计算量更大，通常要降低采样密度，例如把 \texttt{n} 设为 48。

  \begin{block}{环面的两个形状参数}
\begin{Verbatim}[fontsize=\tiny]
def explore_implicit_torus(R=1.55, r=0.45):
    plot_implicit_surface(
        lambda x, y, z: F_torus(x, y, z, R=R, r=r),
        xlim=(-2.4, 2.4), ylim=(-2.4, 2.4), zlim=(-1.2, 1.2),
        n=48,
        title=f"交互实验 2：隐式环面 R={R:.2f}, r={r:.2f}",
    )

interact(
    explore_implicit_torus,
    R=FloatSlider(value=1.55, min=1.0, max=2.0, step=0.05, description="R"),
    r=FloatSlider(value=0.45, min=0.2, max=0.75, step=0.05, description="r"),
);
\end{Verbatim}
  \end{block}
\end{frame}

\begin{frame}{采样密度与性能}
  \begin{tabularx}{\linewidth}{@{}lX@{}}
    \toprule
    参数 & 影响 \\
    \midrule
    参数曲线点数 & 点数越大曲线越平滑，通常几百到一千点足够。 \\
    隐式曲线 \texttt{n} & 二维网格大小约为 \(n^2\)，\texttt{n=301} 或 \texttt{401} 常用于课堂演示。 \\
    参数曲面网格 & 网格大小约为 \(n_u n_v\)，过密会让浏览器交互变慢。 \\
    隐式曲面 \texttt{n} & 三维网格大小约为 \(n^3\)，\texttt{64} 已经是二十多万个采样点。 \\
    \bottomrule
  \end{tabularx}
\end{frame}

\begin{frame}{本专题的实践建议}
  \begin{itemize}
    \item 参数对象优先从参数轴开始设计：一维用 \texttt{t}，二维用 \texttt{u,v}。
    \item 隐式对象优先确定观察窗口和等值 \texttt{level}，再调整采样密度。
    \item 二维讲解优先用 Matplotlib，因为标注、等比例坐标和静态导出更直接。
    \item 三维曲面优先用 Plotly，因为旋转和缩放能帮助判断空间结构。
    \item 三维隐式曲面先确认 \texttt{level} 落在采样值范围内，否则不会有可提取的曲面。
  \end{itemize}
\end{frame}

\end{document}
